
\documentclass{article}
\usepackage{colortbl}
\usepackage{makecell}
\usepackage{multirow}
\usepackage{supertabular}

\begin{document}

\newcounter{utterance}

\centering \large Interaction Transcript for game `cladder', experiment `full\_v1.5\_default', episode 2646 with Qwen3.5{-}27B.
\vspace{24pt}

{ \footnotesize  \setcounter{utterance}{1}
\setlength{\tabcolsep}{0pt}
\begin{supertabular}{c@{$\;$}|p{.15\linewidth}@{}p{.15\linewidth}p{.15\linewidth}p{.15\linewidth}p{.15\linewidth}p{.15\linewidth}}
    \# & \multicolumn{2}{c}{Player} && \multicolumn{2}{c}{Game Master} \\
    \hline

    \theutterance \stepcounter{utterance}  
    & & & \multicolumn{4}{p{0.6\linewidth}}{
        \cellcolor[rgb]{0.9,0.9,0.9}{
            \makecell[{{p{\linewidth}}}]{
                \texttt{\tiny{[P1$\langle$GM]}}
                \texttt{You are an expert in causal inference. The following question is not a typical commonsense query, but rather a meticulously designed question created by a professor specializing in causal inference, intended to assess the students' mastery of the course content.} \\
\texttt{Imagine a self{-}contained, hypothetical world with only the following conditions, and without any unmentioned factors or causal relationships: Gender has a direct effect on smoking and lung cancer. Smoking has a direct effect on tar deposit. Tar deposit has a direct effect on lung cancer. Gender is unobserved. For nonsmokers and with no tar deposit, the probability of lung cancer is 79\%. For nonsmokers and with high tar deposit, the probability of lung cancer is 49\%. For smokers and with no tar deposit, the probability of lung cancer is 71\%. For smokers and with high tar deposit, the probability of lung cancer is 40\%. For nonsmokers, the probability of high tar deposit is 79\%. For smokers, the probability of high tar deposit is 40\%. For smokers, would it be more likely to see lung cancer if the person had been a nonsmoker?} \\
\texttt{Start your answer with 'Yes' or 'No', followed by additional reasoning or evidence to support your explanation.} \\
            }
        }
    }
     \\ \\

    \theutterance \stepcounter{utterance}  
    & \multicolumn{4}{p{0.6\linewidth}}{
        \cellcolor[rgb]{0.561,0.737,0.561}{
            \makecell[{{p{\linewidth}}}]{
                \texttt{\tiny{[P1$\rangle$GM]}}
                \texttt{No} \\
\\ 
\texttt{To determine whether a smoker would be more likely to see lung cancer if they had been a nonsmoker, we must calculate the **Average Treatment Effect (ATE)** of smoking on lung cancer for the population of smokers (specifically, the effect of the counterfactual intervention \$do(S=0)\$ on the outcome \$L\$ for those currently with \$S=1\$).} \\
\\ 
\texttt{In this hypothetical world, the causal structure is defined as:} \\
\texttt{1.  \$G \textbackslash{}to S\$ (Gender affects Smoking)} \\
\texttt{2.  \$G \textbackslash{}to L\$ (Gender affects Lung Cancer)} \\
\texttt{3.  \$S \textbackslash{}to T\$ (Smoking affects Tar)} \\
\texttt{4.  \$T \textbackslash{}to L\$ (Tar affects Lung Cancer)} \\
\\ 
\texttt{Since Gender (\$G\$) is unobserved, we cannot simply condition on it. However, we can calculate the probability of lung cancer under the intervention \$do(S=s)\$ by summing over the mediator (Tar, \$T\$), because the path from \$S\$ to \$L\$ is fully mediated by \$T\$ (given the structure \$S \textbackslash{}to T \textbackslash{}to L\$ and no direct \$S \textbackslash{}to L\$ edge mentioned, though the prompt says "Smoking has a direct effect on tar deposit" and "Tar deposit has a direct effect on lung cancer", implying the path is \$S \textbackslash{}to T \textbackslash{}to L\$. Wait, the prompt says "Gender has a direct effect on smoking and lung cancer". It does *not* explicitly state a direct edge \$S \textbackslash{}to L\$. It states \$S \textbackslash{}to T\$ and \$T \textbackslash{}to L\$. Therefore, the effect of \$S\$ on \$L\$ is entirely mediated by \$T\$.} \\
\\ 
\texttt{The probability of Lung Cancer (\$L=1\$) under the intervention \$do(S=s)\$ is calculated using the **g{-}formula** (or adjustment formula) over the mediator \$T\$:} \\
\texttt{\$\$P(L=1 | do(S=s)) = \textbackslash{}sum\_\{t\} P(L=1 | S=s, T=t) \textbackslash{}cdot P(T=t | S=s)\$\$} \\
\\ 
\texttt{Note: In the g{-}formula for \$do(S=s)\$, the distribution of the mediator \$T\$ is determined by the intervention \$S=s\$, i.e., \$P(T|S=s)\$. The outcome probability \$P(L|S=s, T=t)\$ is the natural conditional probability.} \\
\\ 
\texttt{**Step 1: Calculate \$P(L=1 | do(S=1))\$ (Probability of cancer for a smoker)**} \\
\texttt{We are given:} \\
\texttt{*   \$P(L=1 | S=1, T=0) = 0.71\$} \\
\texttt{*   \$P(L=1 | S=1, T=1) = 0.40\$} \\
\texttt{*   \$P(T=1 | S=1) = 0.40\$ (Probability of high tar} \\
            }
        }
    }
    & & \\ \\

    \theutterance \stepcounter{utterance}  
    & & & \multicolumn{2}{p{0.3\linewidth}}{
        \cellcolor[rgb]{0.9,0.9,0.9}{
            \makecell[{{p{\linewidth}}}]{
                \texttt{\tiny{[GM$|$GM]}}
                \texttt{No} \\
            }
        }
    }
    & & \\ \\

    \theutterance \stepcounter{utterance}  
    & & & \multicolumn{2}{p{0.3\linewidth}}{
        \cellcolor[rgb]{0.9,0.9,0.9}{
            \makecell[{{p{\linewidth}}}]{
                \texttt{\tiny{[GM$|$GM]}}
                \texttt{no} \\
            }
        }
    }
    & & \\ \\

    \theutterance \stepcounter{utterance}  
    & & & \multicolumn{2}{p{0.3\linewidth}}{
        \cellcolor[rgb]{0.9,0.9,0.9}{
            \makecell[{{p{\linewidth}}}]{
                \texttt{\tiny{[GM$|$GM]}}
                \texttt{game\_result = WIN} \\
            }
        }
    }
    & & \\ \\

\end{supertabular}
}

\end{document}
